% Rúbrica: Resolución del problema (cómo el modelo aborda y resuelve el problema) e Interfaz de uso.
% Fuente: backend/despliegue/, backend/modelado/verificar_bandas.py, backend/api/ (main.py, scoring.py y nia/),
% frontend/ y README.md. La estructura del repo y el tag codigo-v3 van en el anexo.
\section{Resolución del problema}\label{sec:resolucion}

\subsection{Arquitectura y despliegue}\label{sec:arquitectura}\label{sec:arquitectura-flujo}\label{sec:arquitectura-build}

La figura~\ref{fig:arquitectura} sigue el camino de los datos. La ingesta bajó las reseñas y las
fichas de Steam y las publicó en releases con tag fijo (sección~\ref{sec:datos-releases}); de ahí
en adelante nada vuelve a consultar Steam. Render construye la API con
\texttt{backend/despliegue/Dockerfile}, y el build falla si algo no cuadra:
\texttt{preparar\_entorno.py} se detiene si el sha256 de data-v3 o de data-v1 no es el publicado,
entrena el modelo con data-v1 y la partición congelada, y \texttt{verificar\_bandas.py} compara la
banda de los \cifraJuegosCatalogo{} juegos con la referencia validada, así que basta con que cambie
una para que falle. El build también rechaza la imagen si le llega una biblioteca de modelos de
texto, que deben calcularse fuera de la API.

\begin{figure}[htbp]
  \centering
  % Generado por documento/generar_figuras.py. No se edita a mano.
\begin{tikzpicture}[caja/.style={draw=#1, fill=#1!7, rounded corners=2pt, align=center, font=\footnotesize,
  minimum width=3.5cm, minimum height=1.25cm, inner sep=3pt}, flecha/.style={-{Stealth[length=5pt]}, draw=neutro, thick}]
\node[caja=neutro] (steam) at (0, 0) {API de Steam\\\texttt{appreviews}\\y \texttt{appdetails}};
\node[caja=neutro] (ingesta) at (3.95, 0) {Ingesta\\\texttt{ingesta\_steam.py}};
\node[caja=neutro] (releases) at (7.9, 0) {Releases con tag fijo\\data-v1 y data-v3\\con su sha256};
\node[caja=neutro] (build) at (11.85, 0) {Build en Render\\entrena con data-v1\\y compara las bandas};
\node[caja=neutro] (notebooks) at (3.95, -1.9) {Notebooks 00, 01 y 02\\código del tag codigo-v6};
\node[caja=acento] (frontend) at (7.9, -1.9) {Frontend Angular\\en Vercel};
\node[caja=acento] (api) at (11.85, -1.9) {API FastAPI\\en Render};
\draw[flecha] (steam) -- (ingesta);
\draw[flecha] (ingesta) -- (releases);
\draw[flecha] (releases) -- (build);
\draw[flecha] (releases.south) -- (notebooks.north east);
\draw[flecha] (build) -- (api);
\draw[flecha, {Stealth[length=5pt]}-{Stealth[length=5pt]}] (frontend) -- (api);
\end{tikzpicture}

  \caption{De la API de Steam a la interfaz. Después de la ingesta, ningún paso consulta Steam.
  Fuente: \texttt{backend/despliegue/Dockerfile}, \texttt{preparar\_entorno.py} y
  \texttt{frontend/vercel.json}.}
  \label{fig:arquitectura}
\end{figure}

La API, hecha con FastAPI, expone \cifraEndpoints{} endpoints: \cifraEndpointsRiesgo{} de riesgo
(\texttt{/catalogo}, \texttt{/perfil}, \texttt{/prediccion}, \texttt{/explicacion} y
\texttt{/panorama}), \cifraEndpointsComunidad{} para calificaciones, comentarios y reacciones, y
\cifraEndpointsNia{} de Nia, descritos en su página \texttt{/docs}\footnote{FastAPI arma esa página
con el contrato OpenAPI de cada endpoint: qué recibe, qué devuelve y qué errores da.}. Todo lo que
entra se valida con Pydantic\footnote{Una biblioteca de Python que revisa el tipo y el rango de cada
dato de entrada y rechaza la solicitud si no cumple.}. \texttt{/prediccion} sigue recibiendo el
perfil por compatibilidad, pero el riesgo que devuelve es del título y no cambia con él. El
frontend, en Angular, es un sitio estático en Vercel.

La API está en el plan gratis de Render: se duerme cuando nadie la usa, y la primera consulta
después tarda en responder. Su disco es efímero, así que las calificaciones, los comentarios y los
votos a Nia se borran al reiniciarse; el catálogo y el modelo vienen dentro de la imagen y no se
pierden. Nia usa un modelo de lenguaje de OpenAI solo si el servidor tiene clave; si no, contesta
por reglas, sin costo. Los comentarios, las reacciones y las consultas y votos a Nia tienen un
límite por minuto, por usuario y por IP.

\subsection{Interfaz de uso}\label{sec:interfaz}
\pendiente{redacción (etapa 3)}

\subsection{Cómo responde la pregunta}\label{sec:resolucion-pregunta}
\pendiente{redacción (etapa 3)}
